
\documentclass[a4paper,11pt]{article}
\usepackage[a4paper, margin=8em]{geometry}

% usa i pacchetti per la scrittura in italiano
\usepackage[french,italian]{babel}
\usepackage[T1]{fontenc}
\usepackage[utf8]{inputenc}
\frenchspacing 

% usa i pacchetti per la formattazione matematica
\usepackage{amsmath, amssymb, amsthm, amsfonts}

% usa altri pacchetti
\usepackage{gensymb}
\usepackage{hyperref}
\usepackage{standalone}

% imposta il titolo
\title{Appunti Comunicazioni Numeriche}
\author{Luca Seggiani}
\date{2026}

% imposta lo stile
% usa helvetica
\usepackage[scaled]{helvet}
% usa palatino
\usepackage{palatino}
% usa un font monospazio guardabile
\usepackage{lmodern}

\renewcommand{\rmdefault}{ppl}
\renewcommand{\sfdefault}{phv}
\renewcommand{\ttdefault}{lmtt}

% disponi il titolo
\makeatletter
\renewcommand{\maketitle} {
	\begin{center} 
		\begin{minipage}[t]{.8\textwidth}
			\textsf{\huge\bfseries \@title} 
		\end{minipage}%
		\begin{minipage}[t]{.2\textwidth}
			\raggedleft \vspace{-1.65em}
			\textsf{\small \@author} \vfill
			\textsf{\small \@date}
		\end{minipage}
		\par
	\end{center}

	\thispagestyle{empty}
	\pagestyle{fancy}
}
\makeatother

% disponi teoremi
\usepackage{tcolorbox}
\newtcolorbox[auto counter, number within=section]{theorem}[2][]{%
	colback=blue!10, 
	colframe=blue!40!black, 
	sharp corners=northwest,
	fonttitle=\sffamily\bfseries, 
	title=~\thetcbcounter: #2, 
	#1
}

% disponi definizioni
\newtcolorbox[auto counter, number within=section]{definition}[2][]{%
	colback=red!10,
	colframe=red!40!black,
	sharp corners=northwest,
	fonttitle=\sffamily\bfseries,
	title=~\thetcbcounter: #2,
	#1
}

% U.D.M
\newcommand{\amp}{\ensuremath{\mathrm{A}}}
\newcommand{\volt}{\ensuremath{\mathrm{V}}}
\newcommand{\meter}{\ensuremath{\mathrm{m}}}
\newcommand{\second}{\ensuremath{\mathrm{s}}}
\newcommand{\farad}{\ensuremath{\mathrm{F}}}
\newcommand{\henry}{\ensuremath{\mathrm{H}}}
\newcommand{\siemens}{\ensuremath{\mathrm{S}}}

% circuiti
\usepackage{circuitikz}
\usetikzlibrary{babel}

% disegni
\usepackage{pgfplots}
\pgfplotsset{width=10cm,compat=1.9}

% disponi codice
\usepackage{listings}
\usepackage[table]{xcolor}

\lstdefinestyle{codestyle}{
		backgroundcolor=\color{black!5}, 
		commentstyle=\color{codegreen},
		keywordstyle=\bfseries\color{magenta},
		numberstyle=\sffamily\tiny\color{black!60},
		stringstyle=\color{green!50!black},
		basicstyle=\ttfamily\footnotesize,
		breakatwhitespace=false,         
		breaklines=true,                 
		captionpos=b,                    
		keepspaces=true,                 
		numbers=left,                    
		numbersep=5pt,                  
		showspaces=false,                
		showstringspaces=false,
		showtabs=false,                  
		tabsize=2
}

\lstdefinestyle{shellstyle}{
		backgroundcolor=\color{black!5}, 
		basicstyle=\ttfamily\footnotesize\color{black}, 
		commentstyle=\color{black}, 
		keywordstyle=\color{black},
		numberstyle=\color{black!5},
		stringstyle=\color{black}, 
		showspaces=false,
		showstringspaces=false, 
		showtabs=false, 
		tabsize=2, 
		numbers=none, 
		breaklines=true
}

\lstdefinelanguage{javascript}{
	keywords={typeof, new, true, false, catch, function, return, null, catch, switch, var, if, in, while, do, else, case, break},
	keywordstyle=\color{blue}\bfseries,
	ndkeywords={class, export, boolean, throw, implements, import, this},
	ndkeywordstyle=\color{darkgray}\bfseries,
	identifierstyle=\color{black},
	sensitive=false,
	comment=[l]{//},
	morecomment=[s]{/*}{*/},
	commentstyle=\color{purple}\ttfamily,
	stringstyle=\color{red}\ttfamily,
	morestring=[b]',
	morestring=[b]"
}

% disponi sezioni
\usepackage{titlesec}

\titleformat{\section}
	{\sffamily\Large\bfseries} 
	{\thesection}{1em}{} 
\titleformat{\subsection}
	{\sffamily\large\bfseries}   
	{\thesubsection}{1em}{} 
\titleformat{\subsubsection}
	{\sffamily\normalsize\bfseries} 
	{\thesubsubsection}{1em}{}

% disponi alberi
\usepackage{forest}

\forestset{
	rectstyle/.style={
		for tree={rectangle,draw,font=\large\sffamily}
	},
	roundstyle/.style={
		for tree={circle,draw,font=\large}
	}
}

% disponi algoritmi
\usepackage{algorithm}
\usepackage{algorithmic}
\makeatletter
\renewcommand{\ALG@name}{Algoritmo}
\makeatother

% disponi numeri di pagina
\usepackage{fancyhdr}
\fancyhf{} 
\fancyfoot[L]{\sffamily{\thepage}}

\makeatletter
\fancyhead[L]{\raisebox{1ex}[0pt][0pt]{\sffamily{\@title \ \@date}}} 
\fancyhead[R]{\raisebox{1ex}[0pt][0pt]{\sffamily{\@author}}}
\makeatother

\begin{document}
% sezione (data)
\section{Lezione del 27-03-26}

% stili pagina
\thispagestyle{empty}
\pagestyle{fancy}

% testo
\subsection{Trasformate continue di Fourier generalizzate}
Iniziamo a vedere le \textbf{TCF} generalizzate, cioè le \textit{Trasformate Continue di Fourier} generalizzate. Una generalizzazione è necessaria in quanto è necessario che i segnali che trasformiamo abbiano energia finita.

Ad esempio, non possiamo calcolare la trasformata di Fourier del gradino di \textit{Heaviside} $u(t)$, in quanto per come è definito:
$$
u(t) =
\begin{cases}
	1, \quad t > 0 \\
	\frac{1}{2}, \quad t = 0 \\
	0, \quad t < 0
\end{cases}
$$
ha energia infinita. La TCF generalizzata vale invece per tutti i segnali ad energia infinita, ad esempio i segnali periodici.

\par\smallskip

\noindent
\textbf{\textsf{Delta di Dirac}}

\noindent
Definiamo innanzitutto la funzione \textbf{delta di Dirac}, $\delta(t)$, come:
$$
\delta(t) = \frac{d}{dt} u(t), \quad u(t) = \int \delta(t) \, dt
$$
già incontrata in fondamenti di automatica. 

Sempre da fondamenti di automatica, ricordiamo che le proprietà importanti della delta di Dirac sono:
\begin{enumerate}
	\item L'\textbf{integrale}:		
		$$
		\int_{-\infty}^{+\infty} \delta(t - t_0) \, dt = 1
		$$
		che diamo per scontato, in quanto non sarebbe in verità possibile trovare una funzione analitica che rispetta tale proprietà (da qui si parla di \textit{distribuzione}, non funzione).
	\item La \textbf{proprietà di campionamento}:	
		$$
		f(t) \delta(t - t_0) = f(t_0) \delta(t - t_0)
		$$
		che si dimostra impostando l'integrale (che assomiglia a una convoluzione ma non vi ricade completamente):
		$$
		\int_{-\infty}^\infty f(t) \delta_\epsilon (t - t_0) \, dt
		$$
		Notiamo subito che l'integrale ha valore $\neq 0$ solo nell'intervallo $[t_0, t_0 + \epsilon]$, e quindi:
		$$
		= \int_{t_0}^{t_0 + \epsilon} \frac{f(t)}{\epsilon} \, dt = \frac{F(t)}{\epsilon} \Bigg|^{t_0 + \epsilon}_{t_0} = \frac{F(t_0 + \epsilon) - F(t_0)}{\epsilon}
		$$
		per una qualche primitiva $F$. Passando al limite $\epsilon \rightarrow 0$, si ha:
		$$
		\lim_{\epsilon \rightarrow 0} \frac{F(t_0 + \epsilon) - F(t_0)}{\epsilon} = f(t_0)
		$$
		Questo significa che vale:
		$$
		\int_{-\infty}^{\infty} f(t) \delta(t - t_0) = \int_{-\infty}^{\infty} f(t_0) \delta(t - t_0) = f(t_0)
		$$
		e visto che la delta di Dirac è definita $\neq 0$ effettivamente in un unico punto in $\mathbb{R}$, dovrà valere la tesi. \qed
	\item Il fatto che la delta di Dirac è \textbf{pari}, chiaramente da come è definita;
	\item Il fatto che la delta di Dirac rappresenta l'\textit{elemento neutro} del \textbf{prodotto di convoluzione}:
		$$
		(f * g)(t) = \int_{-\infty}^{\infty} f(\tau) g(t - \tau) \, d\tau
		$$
		presa infatti $g(t) = \delta(t - t_0)$ si ha:
		$$
		= \int_{-\infty}^{\infty} f(\tau) \delta_\epsilon (\tau - t + t_0) \, d\tau = \int_{t - t_0}^{t - t_0 + \epsilon} \frac{f(\tau)}{\epsilon} \, d\tau = \frac{F(t)}{\epsilon} \Bigg|^{t - t_0 + \epsilon}_{t - t_0} = \frac{F(t - t_0 + \epsilon) - F(t - t_0)}{\epsilon}
		$$
		Passando al limite $\epsilon \rightarrow 0$:
		$$
		\lim_{\epsilon \rightarrow \infty} \frac{F(t - t_0 + \epsilon) - F(t - t_0)}{\epsilon} = f(t - t_0)
		$$
		cioè si ottiene la stessa funzione $f$ scostata di un valore $t_0$, lo stesso di cui è scostata $\delta$, da cui la tesi. \qed
\end{enumerate}

Proviamo quindi a trasformare la trasformata di Fourier della delta di Dirac:
$$
\delta(t) \rightarrow \int_{-\infty}^{\infty} \delta(t) e^{-i 2 \pi f t} \, dt = e^{-i 2 \pi f t} \Big|_{t=0} = 1
$$
direttamente dalla proprietà di campionamento esibita dalla delta di Dirac.

\par\smallskip

Iniziamo quindi a calcolare alcune trasformate di Fourier generalizzate sfruttando le proprietà appena trovate riguardo alla trasformata della delta di Dirac.

\begin{enumerate}
	\item  \textbf{\textsf{Segnale costante}} \\
Prendiamo quindi un segnale ad energia infinita, come ad esempio la funzione costante:
$$
x(t) = 2
$$
A questo punto calcolare la trasformata sarà semplice, in quanto per un segnale DC potremo usare la appena trovata trasformata della delta, assieme al teorema di dualità, per dire:
$$
\rightarrow 2 \cdot \delta(t)
$$

Questo è banale, ma ci dimostra che è effettivamente possibile calcolare trasformate di segnali ad energia infinita. Vedremo adesso come trasformare segnali ad energia infinita più complessi (ma anche più utili);

\item  \textbf{\textsf{Segnale con ritardo}} \\
Nel caso della semplice delta di Dirac con scostamento temporale, si ottiene direttamente dal teorema del ritardo:
$$
\delta(t - t_0) \rightarrow e^{-i 2 \pi f t_0}
$$
da cui segue direttamente (per il teorema della dualità, 7.4):
$$
e^{i 2 \pi f_0 t} \rightarrow \delta(f - f_0)
$$

Questo risultato è estremamente utile in quanto ci mostra come trasformare la forma tipica dalle formule di Eulero:
$$
e^{i 2 \pi f_0 t} = \cos (2 \pi f_0 t) + i \sin (2 \pi f_0 t)
$$

\item \textbf{\textsf{Segnali periodici}} \\
Vediamo quindi come l'ultima proprietà trovata può tornarci utile per calcolare la trasformata di funzioni periodiche. Abbiamo detto che questo non era possibile in quanto è necessario che i segnali che trasformiamo (con la trasformata non generalizzata) abbiano energia finita.

Sfruttando le formule di Eulero, ed applicando la trasformata delta di Dirac, abbiamo invece:
$$
x_c(t) = \cos(2 \pi f_0 t) = \frac{e^{i 2 \pi f_0 t} + e^{-i 2 \pi f_0 t}}{2} \rightarrow
\frac{\delta(f - f_0) + \delta(f + f_0)}{2} = X_c(f)
$$
$$
x_s(t) = \sin(2 \pi f_0 t) = \frac{e^{i 2 \pi f_0 t} - e^{-i 2 \pi f_0 t}}{2} \rightarrow
\frac{\delta(f - f_0) - \delta(f + f_0)}{2j} = X_s(f)
$$

Questo è esattamente in accordo con quanto avevamo visto, seppur anticipandolo in maniera approssimata, introducendo la trasformata di Fourier in 6.1.1 (notiamo che al tempo avevamo usato la DCT, che di per sé mostra questo comportamento della delta di Dirac in maniera "discretizzata").

\par\smallskip

Fissiamo i concetti notando come è possibile dimostrare il teorema della modulazione 7.6 sfruttando il teorema della moltiplicazione e l'appena trovata trasformata del segnale coseno. Avevamo detto che una modulazione non era altro che il prodotto per una cosinusoide: 
$$
z(t) = x(t) \cdot y(t) = x(t) \cos{2 \pi f_0 t}
$$
Questo ci porta ad applicare direttamente la convoluzione nel dominio frequenziale:
$$
Z(f) = X(f) \circledast \frac{\delta(f - f_0) + \delta(f + f_0)}{2}
$$
dalla trasformata del coseno appena calcolata. Ora, visto che la delta di dirac rappresenta il componente neutro della convoluzione (per la proprietà di campionamento), sarà immediato dire:
$$
Z(f) = \frac{X(f - f_0) + X(f + f_0)}{2}
$$
che è esattamente la tesi del teorema di modulazione, che risulta quindi dimostrato. \hfill $\square$

\item \textbf{\textsf{Segnale gradino}} \\
Come ultimio esempio, vediamo come calcolare la trasformata di Fourier generalizzata del classico gradino di \textit{Heaviside}:
$$
u(t) =
\begin{cases}
	0, \quad t < 0 \\
	1, \quad t \geq 0
\end{cases}
$$
In questo caso preferiamo usare la notazione alternativa:
$$
u(t) = \frac{1}{2} + \frac{1}{2} \mathrm{sgn}(t)
$$
Questa ci permette di usare la trasformata (che assumiamo nota, la dimostrazione è nota dal fatto che la funzione segno deriva alla delta di Dirac e dalle proprietà delle trasformate delle derivate, non viste in questi appunti):
$$
\frac{1}{t} \rightarrow -j \pi \mathrm{sgn}(x)
$$
Per cui il calcolo risulta:
$$
u(t) = \frac{1}{2} + \frac{1}{2} \mathrm{sgn}(t) = \frac{1}{2} \delta(t) + \frac{1}{j 2 \pi f}
$$



\end{enumerate}

\end{document}
